\scanpage{386}
\mspage{302}
% archive sheet number: 385

comme \deleted{composés} morphismes de transition.
Cela achève la démonstration de 4.13.10,
et par là de 4.13.9, 4.13.8.

\uline{Remarque} 4.13.10.2. Voici ce qu'on a eu
à utiliser vraiment, pour établir la
conclusion de 4.13.10, en termes de $I,\pi$,
$(J_r)_{r\in A}$ (et sans introduire de $\pi'_0$,\unsure{)}\deleted{\unsure{$M$}}

\textcircled{a'} $I$ est g\textsuperscript{r}de devant tout cardinal \inserted{$\pi'$} $<\pi$.

\textcircled{b'} $I$ est stable par $\varinjlim$ de type $I_{\alpha_\pi}$

\textcircled{c'} $\forall r\in A$, $J_r$ est stable dans $I$ par les
$\varinjlim$ de type $I_{\alpha_\pi}$

\textcircled{\textcircled{d'}} $\forall r\in A$, $J_r$ est cofinal dans $I$

\textcircled{e\unsure{'|''}} $\operatorname{card}A\leq\pi$.

(NB. d') et e') ne sont autres que d) et e)
de 4.3.10, tandis que a'), b'), c'), sont
des versions affaiblies de a), b), c)
respectivement.

\uline{4.13.8} (reformulation) Soit pour $r\in A$, $\pi''_r\geq\pi'_r\geq\pi_0$ (= card. utile \inserted{\unsure{pour $M$}} \unsure{min.})
comme spécifié dans 4.13.2 b) (et c), \deleted{\unsure{et}} soit $\pi'_r=\operatorname{Sup}_{r\in A}\pi'_r$, $\pi''\deleted{\unsure{{}_r}}=\operatorname{Sup}_{r\in A}\pi''_r$, de
sorte que pour tout cardinal $\pi$ où $\pi^{\unsure{\pi'|\pi'_r|\pi'_0}}=\pi\geq\pi''$, $\underline S_r$ soit $\pi$-accessible
dans $M$ (et $\pi$ utile pour $\underline S_r$ \uline{et} pour $M$). Soit $c$ un cardinal \uline{non
limite} tel que \deleted{\unsure{$\operatorname{card} A,\operatorname{Sup}(\pi',\pi'')$}}\inserted{$c\geq\operatorname{card}A$ et $c>\pi'$}
\deleted{\unsure{p.ex.} $\unsure{c=\operatorname{card}A}$ si $\operatorname{card} A$ est \unsure{majoré}}
(p.ex. \deleted{limite, ou} le successeur de \deleted{$\unsure{\operatorname{card}(A)}$}\inserted{$\operatorname{Sup}(\operatorname{card}A,\pi')$}.) Prenons $\pi$ tel
que
\[
\pi^{\pi'}=\pi\geq\operatorname{Sup}(\pi'',c).
\]
Alors $\underline S=\bigcap_{r\in A}\underline S_r$ est $\pi$-accessible.

En effet, si $\alpha_c$ est le premier ordinal de cardinal $c$, comme $c$ n'est
pas limite, $I_{\alpha_c}$ est grand devant tt cardinal $c'<c$, en particulier
devant $\pi'$. \deleted{\unsure{$\geq\pi'$}} Or \deleted{\unsure{les $M_\pi$}} pour tt $x\in\operatorname{Ob}M$, $I=M_{\pi/x}$ et les $J_r=\underline S_{r/x}$
sont stables par $\varinjlim$ g\textsuperscript{r}des dev\textsuperscript{t} $\pi'$ et de cardinal $\leq\pi$, \deleted{\unsure{donc} \unsure{les}}
